% docs/physics/latex/06_kaon.tex - SIMC++ physics note 06: kaon electroproduction N(e,e'K+)Y.
% Standalone PDF docs/physics/pdf/SIMCpp_Physics_06_kaon.pdf (copied to src/reactions/kaon/ and src/physics/kaon/);
% a chapter of the book SIMC++ Physics Reference. Rules: docs/physics/README.md; structure: template_reaction.tex.
\documentclass{simcnote}
\simccommon
\begin{document}
\simcnotemeta{06}{Kaon electroproduction}{Kaon electroproduction: H, D, A(e,e\pr K\texorpdfstring{$^+$}{+})Y and coherent hypernuclei}%
  {Reaction plug-in \texttt{kaon\_production} with the physics model \texttt{kaon}: the factorized cross section, the
   Saghai amplitude tables, \texorpdfstring{$\Lambda$, $\Sigma^0$, $\Sigma^-$}{Lambda, Sigma0, Sigma-} and bound final states, and the kaon decay in flight}
\simcnotemaketitle
\simcabstract{Exclusive K$^+$ electroproduction with the electron and the kaon detected in coincidence and the hyperon
($\Lambda$, $\Sigma^0$ or $\Sigma^-$) or the hypernucleus missing. Both directions and the electron energy are
generated; the kaon momentum follows from two-body kinematics on a nucleon that moves with a Fermi momentum for D and
heavier targets. The weight is a factorized fit in $Q^2$, $W$ and $t$ boosted from the photon--nucleon centre of mass,
times the virtual-photon flux; the Saghai amplitude model is evaluated for a diagnostic column only. The kaon either
decays in flight in the spectrometer or the weight carries its survival probability. Notation and common physics: note
00.}
\simcnotetoc

% =====================================================================================================================
\section{The reaction in SIMC++}\lbl{sec:reaction}

\subsection{Process and scope}
The reaction is $e + N \to e' + K^+ + Y$ with $Y=\Lambda$ or $\Sigma^0$ on the proton and $Y=\Sigma^-$ on the neutron,
on hydrogen (quasi-free on a proton at rest), deuterium and heavier nuclei (quasi-free on a moving nucleon), or the
coherent production $e + A \to e' + K^+ + {}_Y\!A'$ of a bound hypernucleus, treated as production on a ``heavy
nucleon'' of mass $M_A$. The kaon is detected in the hadron arm; the hyperon is the missing particle and is not
followed (it neither decays nor radiates).

What is simulated: the beam, the vertex, energy loss and multiple scattering, radiation on the three legs (the kaon is
the radiating hadron), the Fermi motion of the struck nucleon (D and $A\ge3$), the spectrometers, the kaon decay in
flight inside the spectrometers (\code{doing\_decay} = 1) or, without it, a survival weight, and the factorized cross
section of \secr{sec:xsec}. What is not simulated: final-state interactions and nuclear transparency, Pauli blocking,
the decay of the hyperon, any $\varphi$ dependence of the cross section (the factorized model has none, the Saghai
diagnostic is evaluated at $\varphi=0$), and a nuclear model for coherent production (the nucleon fit is used with the
photon--nucleus $W$).

\subsection{Components}
\begin{simctable}{footnotesize}{L{0.15\linewidth}L{0.15\linewidth}L{0.16\linewidth}L{0.25\linewidth}L{0.29\linewidth}}
\caption{The SIMC++ components of kaon production (one component per row).}\lbl{tab:components}\\
\toprule\rowcolor{simchead}\simcth{role} & \simcth{interface} & \simcth{registered name} & \simcth{class} & \simcth{file}\\\midrule\endfirsthead
\toprule\rowcolor{simchead}\simcth{role} & \simcth{interface} & \simcth{registered name} & \simcth{class} & \simcth{file}\\\midrule\endhead
\bottomrule\endlastfoot
reaction & \code{IReaction} & \code{kaon\_production} & \code{KaonProduction} & \file{reactions/kaon/KaonProduction.cpp}\\
generation and vertex completion & (base class) & -- & \code{PionProduction} & \file{reactions/pion/PionProduction.cpp}\\
reconstruction, generation volume & (base class) & -- & \code{EepReactionBase} & \file{reactions/EepReactionBase.cpp}\\
centre-of-mass transformation & (function) & -- & \code{reactions::transformToCm} & \file{reactions/EventChain.cpp}\\
cross section & \code{IPhysicsModel} & \code{kaon} & \code{kaon::KaonModel} & \file{physics/kaon/KaonModel.cpp}\\
cross-section kernel & (function) & -- & \code{kaon::peeK} & \file{physics/kaon/KaonCrossSections.cpp}\\
factorized model & (function) & -- & \code{kaon::sigFactorized} & \file{physics/kaon/KaonCrossSections.cpp}\\
Saghai model, $\Lambda$ & (function) & -- & \code{kaon::eekeek} & \file{physics/kaon/KaonCrossSections.cpp}\\
Saghai model, $\Sigma$ & (function) & -- & \code{kaon::eekeeks} & \file{physics/kaon/KaonCrossSections.cpp}\\
table interpolation & (function) & -- & \code{kaon::fint} & \file{physics/kaon/KaonCrossSections.cpp}\\
amplitude tables & (class) & -- & \code{kaon::SaghaiTables} & \file{physics/kaon/KaonCrossSections.cpp} (files in \file{data/kaon})\\
Fermi momenta & (function) & -- & \code{ReadMomentumDistribution} & \file{setup/SetupTheory.cpp}\\
missing energy for $A\ge3$ & (class) & -- & \code{nuclear::SfTable::generateEm} & \file{physics/nuclear/SpectralFunctionTables.cpp}\\
decay in flight & (functions) & -- & \code{decayProducts}, \code{project} & \file{optics/CosyEngine.cpp} (used by the spectrometer models)\\
radiation & \code{IRadiator} & \code{simc\_rad} & \code{SimcRadiator} & \file{radiation/SimcRadiator.cpp}\\
spectrometers & \code{ISpectrometer} & from the arm codes & e.g.\ \code{shms::ShmsSpectrometer} & spectrometers/\\
\end{simctable}

\subsection{Selection and override}
\code{doing\_kaon} = 1 selects \code{kaon\_production} with the model \code{kaon} (\ctabr{tab:defaults}), unless
\code{doing\_semi} is also set (then the run is kaon SIDIS, another note) or \code{doing\_phsp} = 1. The set-up
(\code{DeriveReactionFlags}, \code{BuildReactionAndKinematics}) sets $M_h = M_K = 493.677$ MeV and the masses of
\tabr{tab:masses}; the target flags are ``hydrogen'' for $A=1$ and for every coherent case, ``deuterium'' for $A=2$ and
``nucleus'' for $A\ge3$ \fsrc{dbase.f:163-180}. $\Sigma^-$ on hydrogen, coherent production on H or D and any other
\code{which\_kaon} stop the run at start-up. For coherent production the input final-nucleus mass is replaced twice
(first by $M_A-M_p$ as for every heavy target, then by the guess of the table) and is therefore never used.

\begin{simctable}{footnotesize}{L{0.11\linewidth}L{0.23\linewidth}L{0.16\linewidth}L{0.2\linewidth}L{0.3\linewidth}}
\caption{Final states (\code{which\_kaon}; one value per row; $n_\mathrm{tgt}$ is the nucleon count of the weight).}\lbl{tab:masses}\\
\toprule\rowcolor{simchead}\simcth{which\_kaon} & \simcth{process} & \simcth{$M_N$ (struck)} & \simcth{$M_Y$ (recoil)} & \simcth{$n_\mathrm{tgt}$}\\\midrule\endfirsthead
\toprule\rowcolor{simchead}\simcth{which\_kaon} & \simcth{process} & \simcth{$M_N$ (struck)} & \simcth{$M_Y$ (recoil)} & \simcth{$n_\mathrm{tgt}$}\\\midrule\endhead
\bottomrule\endlastfoot
0 & $p(e,e'K^+)\Lambda$ & $M_p$ & $M_\Lambda = 1115.68$ MeV & $Z$\\
1 & $p(e,e'K^+)\Sigma^0$ & $M_p$ & $M_{\Sigma^0} = 1192.64$ MeV & $Z$\\
2 & $n(e,e'K^+)\Sigma^-$ & $M_n$ & $M_{\Sigma^-} = 1197.45$ MeV & $N$\\
10 & $A(e,e'K^+){}_\Lambda A'$ & $M_A$ & $M_A - M_p + M_\Lambda$ & $Z$\\
11 & $A(e,e'K^+){}_{\Sigma^0} A'$ & $M_A$ & $M_A - M_p + M_{\Sigma^0}$ & $Z$\\
12 & $A(e,e'K^+){}_{\Sigma^-} A'$ & $M_A$ & $M_A - M_n + M_{\Sigma^-}$ & $N$\\
\end{simctable}
In this note $M_N$ is the struck mass and $M_Y$ the recoil mass of \tabr{tab:masses}; for coherent production there is
no residual system ($M_\mathrm{rec}=0$).

The model reads its four tables from the data directory's \file{kaon/} (\code{l1\_kaon\_data\_files} checks that they
are the Fortran files). The JSON input can override the model or its options:
\begin{simctable}{footnotesize}{L{0.27\linewidth}L{0.3\linewidth}L{0.37\linewidth}}
\caption{JSON keys of the reaction and the model (one key per row).}\lbl{tab:json}\\
\toprule\rowcolor{simchead}\simcth{key} & \simcth{meaning} & \simcth{example}\\\midrule\endfirsthead
\toprule\rowcolor{simchead}\simcth{key} & \simcth{meaning} & \simcth{example}\\\midrule\endhead
\bottomrule\endlastfoot
\code{pipeline.reaction} & the reaction & \code{"kaon\_production"}\\
\code{pipeline.physics.type} & the cross-section model & \code{"kaon"}\\
\code{pipeline.physics.data\_dir} & directory of the four tables & \code{"data/kaon"}\\
\code{pipeline.physics.verbose} & print the table messages and the model warnings & \code{true}\\
\end{simctable}
A replacement model must return the five-fold cross section of \secr{sec:xsec} in $\mu$b/(MeV\,sr$^2$), and, when the
decay is not simulated, provide the survival probability in the event (\code{main.ntup.decdist}) as \code{kaon} does.

\subsection{Event flow}
\begin{simctable}{footnotesize}{L{0.05\linewidth}L{0.33\linewidth}L{0.62\linewidth}}
\caption{One try of kaon production (one step per row; the general chain is \ctabr{tab:chain}).}\lbl{tab:flow}\\
\toprule\rowcolor{simchead}\simcth{\#} & \simcth{C++ call} & \simcth{what is specific to this reaction}\\\midrule\endfirsthead
\toprule\rowcolor{simchead}\simcth{\#} & \simcth{C++ call} & \simcth{what is specific to this reaction}\\\midrule\endhead
\bottomrule\endlastfoot
0 & \code{PionProduction::prepareCentral} & once per run: the central-event quantities read by the model (\secr{sec:central})\\
1 & \code{StandardBeam::generate} & nothing\\
2 & \code{PionProduction::generate} & both pairs of slopes, $E'$, and for D/A the Fermi momentum and $E_m$ (\eqsr{eq:genslopes}{eq:fermi})\\
3 & \code{PionProduction::completeEv} & the kaon energy from the two-body equation (\eqr{eq:quad}), $W$, $\varepsilon$, $t$, $\varphi_{pq}$, both solid-angle factors\\
4 & \code{SimcRadiator::radiate} & photon limits of \eqr{eq:radlimits}; a tail-1 photon recomputes the kaon at the lower beam energy\\
5 & \code{MonteCarlo::run} & the kaon crosses the hadron arm; with \code{doing\_decay} it may decay in flight (\eqr{eq:decay})\\
6 & \code{EepReactionBase::completeRecon} & missing mass and $-t$ of the reconstructed event (\eqr{eq:recon})\\
7 & \code{KaonProduction::completeMain} & \code{KaonModel::sigccMain} (\eqr{eq:peek}), $n_\mathrm{tgt}$, the survival probability (\eqr{eq:weight})\\
\end{simctable}

% =====================================================================================================================
\section{Kinematics, generated variables and generation volume}\lbl{sec:gen}
\subsection{Generated variables}
After the beam part (\csecr{sec:beam}) the reaction draws, with uniform random numbers $\xi_i$:
\begin{feqa}{reactions/pion/PionProduction.cpp: PionProduction::generate}{event.f:276-316}
y'_e &= y'_{e,\mathrm{min}} + \xi_1\Delta y'_e,\qquad x'_e = x'_{e,\mathrm{min}} + \xi_2\Delta x'_e,\nonumber\\
y'_h &= y'_{h,\mathrm{min}} + \xi_3\Delta y'_h,\qquad x'_h = x'_{h,\mathrm{min}} + \xi_4\Delta x'_h,\nonumber\\
E' &= E'_\mathrm{lo} + \xi_5\,(E'_\mathrm{hi}-E'_\mathrm{lo}),\nonumber\\
E'_\mathrm{lo} &= \max(E'^\mathrm{gen}_\mathrm{min},\Sigma_\mathrm{min}),\qquad E'_\mathrm{hi} = \min(E'^\mathrm{gen}_\mathrm{max},\Sigma_\mathrm{max}),
\lbl{eq:genslopes}
\end{feqa}
with the generation weight of \ceqr{eq:genweight} for $E'$ (1 in practice, since the generation range of $E'$ already lies
inside $[\Sigma_\mathrm{min},\Sigma_\mathrm{max}]$); the try fails if $E'_\mathrm{lo}>E'_\mathrm{hi}$. Both directions are
converted to physics angles (\ceqr{eq:physang}). The kaon momentum is not generated.

\paragraph{Fermi motion (D and $A\ge3$, quasi-free only).} The struck nucleon gets a momentum $p_F$ drawn from the
cumulative momentum distribution $C(p_i)$ of the target (\code{deut.dat}; \code{he3.dat}, \code{he4.dat},
\code{c12.dat} for $A=3,4,12$; \code{c12.dat} with a message for other $A$), normalised to $C(p_n)=1$, and an isotropic
direction $\hat p_F$:
\begin{feqa}{reactions/pion/PionProduction.cpp: PionProduction::generate}{event.f:336-372; dbase.f:567-591}
i &= \min\{i:\ \xi_6\le C(p_i)\}\ (\le n),\qquad p_F = p_\mathrm{lo}+\xi_7\,(p_\mathrm{hi}-p_\mathrm{lo}),\nonumber\\
p_\mathrm{lo} &= \tfrac12(p_{i-1}+p_i)\ (0\text{ for }i=1),\qquad p_\mathrm{hi} = \tfrac12(p_i+p_{i+1})\ (p_n\text{ for }i=n),\nonumber\\
\cos\theta_F &= 2\xi_8-1,\qquad \varphi_F = 2\pi\xi_9,\qquad
\hat p_F = (\sin\theta_F\cos\varphi_F,\ \sin\theta_F\sin\varphi_F,\ \cos\theta_F),\nonumber\\
E_m &= M_p+M_n-M_d\ (\text{D}),\qquad E_m = E_m^\mathrm{SF}(p_F)\ (A\ge3),\nonumber\\
E_F &= M_A - \sqrt{(M_A - M_N + E_m)^2 + p_F^2},
\lbl{eq:fermi}
\end{feqa}
where $E_m^\mathrm{SF}(p_F)$ is drawn from the Benhar spectral function of the target at $P_m=p_F$ (the $E_m$
distribution summed over the table nodes, a node chosen with one uniform number and smeared flat over its bin with a
second; the neutron columns for $\Sigma^-$), and $E_F$ is the energy of the off-shell struck nucleon when the
spectator system of mass $M_A-M_N+E_m$ is on shell with momentum $-\vec p_F$. For hydrogen and coherent production
$p_F=0$, $E_m=0$ and $E_F=M_N$.

\subsection{Limits}
The vertex limits of the missing energy and momentum of the spectator system are
\begin{feqa}{setup/SetupLimits.cpp: BuildLimits}{init.f:344-358, 374-402}
\text{H, coherent:}&\quad E_m = P_m = M^\mathrm{V}_\mathrm{rec} = T^\mathrm{V}_\mathrm{rec} = 0,\nonumber\\
\text{D:}&\quad E_m = M_p+M_n-M_d,\qquad P_m\in[0,\ p_n],\nonumber\\
A\ge3:&\quad E_m\in[M_N+M_\mathrm{rec}-M_A,\ E_m^\mathrm{SF,max}],\qquad P_m\in[0,\ p_n],\nonumber\\
&\quad M^\mathrm{V}_\mathrm{rec} = M_A-M_N+E_m,\qquad
T_\mathrm{rec}^\mathrm{min} = \sqrt{(M^\mathrm{V,max}_\mathrm{rec})^2+P_{m,\mathrm{min}}^2}-M^\mathrm{V,max}_\mathrm{rec},\nonumber\\
&\quad T_\mathrm{rec}^\mathrm{max} = \sqrt{(M^\mathrm{V,min}_\mathrm{rec})^2+P_{m,\mathrm{max}}^2}-M^\mathrm{V,min}_\mathrm{rec},\nonumber\\
&\quad T_Y\in\bigl[0,\ E_b^\mathrm{max}+M_N-M_Y-E'_\mathrm{min}-E_{h,\mathrm{min}}-E_{m,\mathrm{min}}-T_\mathrm{rec}^\mathrm{min}\bigr],
\lbl{eq:vertexlimits}
\end{feqa}
with $p_n$ the last node of the momentum distribution, $E_m^\mathrm{SF,max}$ the last $E_m$ node of the spectral
function and $T_Y$ the kinetic energy of the recoiling hyperon, for which only energy conservation gives a limit.
The largest photon energy and the range of the generated electron energy then follow from energy conservation:
\begin{feqa}{setup/SetupLimits.cpp: BuildLimits}{init.f:411-425, 456-469, 477-488}
E_\gamma^\mathrm{tot,max} &= E_b^\mathrm{max}+M_N-M_Y-E'_\mathrm{min}-E_{h,\mathrm{min}}-E_{m,\mathrm{min}}-T^\mathrm{min}_\mathrm{rec},\nonumber\\
\Sigma_\mathrm{max} &= \min\bigl(E_b^\mathrm{max}+M_N-M_Y-E_{h,\mathrm{min}}-E_{m,\mathrm{min}}-T^\mathrm{min}_\mathrm{rec},\ E'_\mathrm{max}+E_{\gamma2}^\mathrm{max}\bigr)+d_E,\nonumber\\
\Sigma_\mathrm{min} &= \max\bigl(E_b^\mathrm{min}+M_N-M_Y-E_{h,\mathrm{max}}-E_{m,\mathrm{max}}-T^\mathrm{max}_\mathrm{rec}-T_Y^\mathrm{max}-E_\gamma^\mathrm{tot,max},\nonumber\\
&\qquad\quad E'_\mathrm{min}\bigr)-d_E\ \ (\ge0),\nonumber\\
[E'^\mathrm{gen}_\mathrm{min},E'^\mathrm{gen}_\mathrm{max}] &= [\max(\Sigma_\mathrm{min},E'_\mathrm{min}),\ \min(\Sigma_\mathrm{max},E'_\mathrm{max}+E_{\gamma2}^\mathrm{max})],
\lbl{eq:sumegen}
\end{feqa}
($E_\gamma^\mathrm{tot,max}=0$ without radiation, \code{Egamma\_gen\_max} when that is set; each enabled tail has
$E_{\gamma i}^\mathrm{max}=E_\gamma^\mathrm{tot,max}$), where $\Sigma$ is here the generated electron energy. The hadron
energy range $[E_{h,\mathrm{min}},E_{h,\mathrm{max}}+E^\mathrm{max}_{\gamma3}]$ serves only the diagnostic histograms;
the kaon slopes are generated inside the hadron edges.

\subsection{Generation volume}
\begin{feq}{reactions/EepReactionBase.cpp: EepReactionBase::genVolume}{simc.f:383-397}
V_\mathrm{gen} = \Delta\Omega_e\,\Delta\Omega_h\,(E'^\mathrm{gen}_\mathrm{max}-E'^\mathrm{gen}_\mathrm{min})\quad[\mathrm{sr^2\,MeV}].
\lbl{eq:genvol}
\end{feq}
The Fermi momentum and $E_m$ are drawn from normalised distributions and add no factor to $V_\mathrm{gen}$.

% =====================================================================================================================
\section{Completing the event}\lbl{sec:complete}
\subsection{Vertex}
With $\nu$, $Q^2$, $\vec q$ of \ceqr{eq:nuq2} and the kaon direction $\hat p$, the kaon energy is the solution of
$\nu + E_F = E_h + \sqrt{M_Y^2+|\vec q+\vec p_F-\vec p_h|^2}$ at fixed directions:
\begin{feqa}{reactions/pion/PionProduction.cpp: PionProduction::completeEv}{event.f:634-698}
a &= -|\vec q|\,\hat q\cdot\hat p\ \bigl[-\,p_F\,\hat p_F\cdot\hat p\bigr],\qquad
b = |\vec q|^2\ \bigl[+\,p_F^2+2|\vec q|p_F\,\hat p_F\cdot\hat q\bigr],\qquad
c = \nu + M_A\ \ \bigl[\nu+E_F\bigr],\nonumber\\
k &= c^2-b+M_h^2-M_Y^2,\qquad Q_A = 4(a^2-c^2),\qquad Q_B = 4ck,\qquad Q_C = -4a^2M_h^2-k^2,\nonumber\\
E_h &= \frac{-Q_B-\sqrt{Q_B^2-4Q_AQ_C}}{2Q_A},\qquad P_h = \sqrt{E_h^2-M_h^2},\qquad \delta_h = 100\,(P_h-P_h^0)/P_h^0,
\lbl{eq:quad}
\end{feqa}
the brackets applying to D and $A\ge3$ ($M_A=M_p$ for hydrogen and the nucleus mass for coherent production, where
$E_F=M_A$ is the same thing). The quadratic is the square of $k-2cE_h=2aP_h$; since $Q_A<0$ the chosen root is the
larger one, the kaon going forward in the centre of mass. The try fails when the discriminant is negative, when
$E_h<0$, when the recoil energy $c-E_h\le M_Y$ or when $E_h\le M_h$. When the other root falls inside the hadron
energy edges a warning is printed (only in the first 5000 tries): the backward kaon is assumed never to be accepted.
The meson variables of the vertex follow (\ceqsr{eq:w}{eq:epsilon}, \ceqr{eq:phipq}):
\begin{feqa}{reactions/pion/PionProduction.cpp: PionProduction::pionMainKinematics}{event.f:723-771}
W^2 &= M_N^2+2M_N\nu-Q^2,\qquad \theta_{pq} = \arccos(\hat p\cdot\hat q),\nonumber\\
-t &= Q^2-M_h^2+2\nu E_h-2P_h|\vec q|\cos\theta_{pq},\qquad
-t_\mathrm{min} = Q^2-M_h^2+2\nu E_h-2P_h|\vec q|,
\lbl{eq:mesonvars}
\end{feqa}
$W$ being the photon--nucleon mass for a nucleon at rest (the photon--nucleus mass for coherent production; the model
uses its own $W$ with the Fermi momentum, \eqr{eq:wcm}). $-t$ is positive; the code stores $-t$ under the name $t$.
The spectator system has $P_m=p_F$, $M_\mathrm{rec}=M_A-M_N+E_m$ and
\begin{feq}{reactions/pion/PionProduction.cpp: PionProduction::completeEv}{event.f:634-635, 947-951}
T_\mathrm{rec} = \sqrt{M_\mathrm{rec}^2+p_F^2}-M_\mathrm{rec}\quad(\text{D},\ A\ge3),\qquad T_\mathrm{rec}=0\quad(\text{H, coherent}).
\lbl{eq:trec}
\end{feq}
The vertex $\vec P_m$, $E_\mathrm{miss}$ (\ceqr{eq:pm}), $z$ and $p_T^2$ are filled as for every reaction. The kaon is
not limited by the generation: its momentum may lie anywhere, and the hadron-arm model decides.

\subsection{Reconstruction}
From the reconstructed electron and kaon (\ceqr{eq:recon}):
\begin{feqa}{reactions/EepReactionBase.cpp: EepReactionBase::completeRecon}{event.f:1325-1336, 1353-1354}
E_m^\mathrm{rec} &= \nu^\mathrm{rec}+M_N-E_h^\mathrm{rec},\qquad
M_x = \operatorname{sign}(m^2)\sqrt{|m^2|},\ \ m^2 = (E_m^\mathrm{rec})^2-(P_m^\mathrm{rec})^2,\nonumber\\
M_x^A &= \operatorname{sign}(m_A^2)\sqrt{|m_A^2|},\ \ m_A^2 = (\nu^\mathrm{rec}+M_A-E_h^\mathrm{rec})^2-(P_m^\mathrm{rec})^2,\nonumber\\
-t^\mathrm{rec} &= (Q^\mathrm{rec})^2-M_h^2+2\bigl(\nu^\mathrm{rec}E_h^\mathrm{rec}-P_h^\mathrm{rec}|\vec q^\mathrm{rec}|\cos\theta^\mathrm{rec}_{pq}\bigr),
\lbl{eq:recon}
\end{feqa}
$M_x$ (tree column \code{missmass}) peaks at $M_Y$ for hydrogen; $M_x^A$ (\code{mmnuc}) is the missing mass of the
whole target. The reconstructed $W$ uses $M_p$ (\ceqr{eq:w}).

% =====================================================================================================================
\section{Jacobians}\lbl{sec:jac}
Both pairs of slopes are generated flat, so the event carries both solid-angle factors (\ceqr{eq:jacomega}):
\begin{feq}{reactions/pion/PionProduction.cpp: PionProduction::completeEv}{event.f:1015-1025}
J = \bigl(1+x_e'^2+y_e'^2\bigr)^{-3/2}\bigl(1+x_h'^2+y_h'^2\bigr)^{-3/2}.
\lbl{eq:jac}
\end{feq}
The transformation of the cross section from the centre of mass to the laboratory solid angle of the kaon is part of the
model (\eqr{eq:jcm}), not of $J$.

% =====================================================================================================================
\section{Cross-section model}\lbl{sec:xsec}
\subsection{What the reaction requires from a model}
\code{KaonProduction::completeMain} sets $S=1$ and calls \code{sigccMain(vertex, main)}, which must return the
five-fold cross section $\mathrm{d}^5\sigma/\mathrm{d}E'\mathrm{d}\Omega_e\mathrm{d}\Omega_h$ in $\mu$b/(MeV\,sr$^2$)
in the laboratory and may write the centre-of-mass quantities and the diagnostic cross sections into the event
(tree columns \code{sigcm}, \code{saghai}, \code{factor}). When \code{doing\_decay} = 0 the model must also store the
survival probability in \code{main.ntup.decdist}; the reaction multiplies the weight by it. The model reads the vertex
$E_\mathrm{in}$, $E'$, $\nu$, $Q^2$, $\vec q$, the kaon four-momentum, $\varepsilon$, $\theta_{pq}$, $\varphi_{pq}$,
$-t$, the Fermi momentum, and the focal-plane path and slopes of the kaon.

\subsection{The model \code{kaon}: centre-of-mass frame}
The photon--nucleon centre of mass moves with
\begin{feqa}{reactions/EventChain.cpp: transformToCm}{jacobians.f:57-133}
\vec\beta^* &= \frac{|\vec q|\,\hat q+p_F\,\hat p_F}{\nu+E_F},\qquad \gamma^* = (1-\beta^{*2})^{-1/2},\nonumber\\
(\nu^*,\vec q^{\,*}) &= \Lambda(\vec\beta^*)\,(\nu,\vec q),\qquad (E^*_h,\vec p^{\,*}_h) = \Lambda(\vec\beta^*)\,(E_h,\vec p_h),\qquad
\theta^* = \angle(\vec p^{\,*}_h,\vec q^{\,*}),
\lbl{eq:cm}
\end{feqa}
$\Lambda$ being the Lorentz boost of \code{loren} written in the frame with $z\parallel\hat q$ and $x$ in the scattering
plane (the frame of $\varphi_{pq}$). The azimuth $\varphi^*$ of $\vec p^{\,*}_h$ is measured around $\vec q^{\,*}$ from
the plane of $\vec q^{\,*}$ and the boosted beam. The invariant mass with the Fermi momentum is
\begin{feq}{physics/kaon/KaonCrossSections.cpp: peeK}{physics\_kaon.f:66-71}
W^2_\mathrm{cm} = s = (\nu+E_F)^2-(|\vec q|+p_F\cos\theta_{Fq})^2-(p_F\sin\theta_{Fq})^2,\qquad \cos\theta_{Fq} = \hat p_F\cdot\hat q.
\lbl{eq:wcm}
\end{feq}
The jacobian from $(t,\varphi^*)$ to the laboratory variables $(\cos\theta_{pq},\varphi_{pq})$ is computed in full
for a non-collinear boost, and divided by $\mathrm{d}t/\mathrm{d}\cos\theta^* = 2P_h^*|\vec q^{\,*}|$ (Mandelstam $t$):
\begin{feqa}{reactions/EventChain.cpp: transformToCm; physics/kaon/KaonCrossSections.cpp: peeK}{jacobians.f:208-271; physics\_kaon.f:52}
J_\mathrm{cm} &= \frac{1}{2P_h^*|\vec q^{\,*}|}\left|\frac{\partial t}{\partial\cos\theta_{pq}}\frac{\partial\varphi^*}{\partial\varphi_{pq}}
 - \frac{\partial t}{\partial\varphi_{pq}}\frac{\partial\varphi^*}{\partial\cos\theta_{pq}}\right|
 = \frac{\mathrm{d}\Omega^*}{\mathrm{d}\Omega_h},\nonumber\\
\frac{\partial t}{\partial\cos\theta_{pq}} &= 2\Bigl[|\vec q|P_h+\Bigl(|\vec q|\cos\theta_{pq}-\frac{\nu P_h}{E_h}\Bigr)\frac{\partial P_h}{\partial\cos\theta_{pq}}\Bigr],\qquad
\frac{\partial t}{\partial\varphi_{pq}} = 2\Bigl(|\vec q|\cos\theta_{pq}-\frac{\nu P_h}{E_h}\Bigr)\frac{\partial P_h}{\partial\varphi_{pq}},
\lbl{eq:jcm}
\end{feqa}
with $\partial P_h/\partial\cos\theta_{pq}$ and $\partial P_h/\partial\varphi_{pq}$ from energy conservation at fixed
$\vec p_F$, and the derivatives of $\varphi^*$ from the boost (\fsrc{jacobians.f:208-269}, transcribed term by term).
Without Fermi motion the boost is along $\vec q$, $\partial\varphi^*/\partial\varphi_{pq}=1$ and the second term
vanishes. The older approximation for a collinear boost is computed too (tree diagnostic only, not in the weight):
\begin{feq}{reactions/EventChain.cpp: transformToCm}{jacobians.f:275-278}
J_\mathrm{cm}^\mathrm{coll}\cdot 2P_h^*|\vec q^{\,*}| = \frac{2\,(E_F-2p_{F}\hat p_{F,z}\frac{E_h}{P_h}\cos\theta_{pq})(|\vec q|+p_F\hat p_{F,z})P_h}
{E_F+\nu-(|\vec q|+p_F\hat p_{F,z})\frac{E_h}{P_h}\cos\theta_{pq}}-2P_hp_F.
\lbl{eq:jold}
\end{feq}

\subsection{The model \code{kaon}: the cross section}
\begin{feqa}{physics/kaon/KaonCrossSections.cpp: peeK}{physics\_kaon.f:52, 96, 108, 132-135}
\sigma &= \sigma_\mathrm{fact}\cdot J_\mathrm{cm}\cdot\Gamma'\cdot f\quad[\mu\mathrm{b/(MeV\,sr^2)}],\nonumber\\
\Gamma' &= \frac{\alpha}{2\pi^2}\,\frac{E'}{E_\mathrm{in}}\,\frac{k_\gamma}{Q^2}\,\frac{1}{1-\varepsilon},\qquad
k_\gamma = \frac{W_\mathrm{cm}^2-M_N^2}{2M_N},\qquad
f = \frac{M_N}{E_F-p_F\hat p_{F,z}},
\lbl{eq:peek}
\end{feqa}
where $\sigma_\mathrm{fact}$ is the centre-of-mass cross section $\mathrm{d}\sigma/\mathrm{d}\Omega^*$ in $\mu$b/sr
(code comment \fsrc{physics\_kaon.f:111-116}), $\Gamma'$ the virtual-photon flux per MeV and sr with the equivalent
photon energy $k_\gamma$ of Hand, and $f$ the correction of the flux for a moving target: the relative velocity
$1-p_F\hat p_{F,z}/E_F$ along the beam and the target energy $E_F$ in place of $M_N$ (code comment
\fsrc{physics\_kaon.f:118-129}; $f=1$ without Fermi motion). The code comments record that $k_\gamma$ with $M_N$
(rather than $E_F-p_F\cos\theta_{Fq}$) was chosen to match the phase space of the generation.

\paragraph{Factorized model (the one in the weight).} The fit of D.~Koltenuk's thesis \citer{koltenuk}: at fixed $Q^2$
and $W$ the cross section falls as $\exp(-b|t-t_\mathrm{min}|)$. The routine receives $Q^2$, $W_\mathrm{cm}$, $-t$ and
$P_h^*$ (all in MeV), recomputes the minimum $-t$ for a proton at rest, and evaluates (in GeV):
\begin{feqa}{physics/kaon/KaonCrossSections.cpp: sigFactorized}{physics\_kaon.f:194-234}
\nu_W &= \frac{W^2+Q^2-M_p^2}{2M_p},\qquad q^*_W = \sqrt{Q^2+\nu_W^2}\,\frac{M_p}{W},\qquad \nu^*_W = \sqrt{q^{*2}_W-Q^2},\nonumber\\
-t_\mathrm{min} &= Q^2-M_K^2+2\nu^*_W\sqrt{P_h^{*2}+M_K^2}-2q^*_WP^*_h,\nonumber\\
F_Q &= (Q^2+A_Q)^{-2},\qquad F_t = \exp\bigl[-b\,\bigl((-t)-(-t_\mathrm{min})\bigr)\bigr],\nonumber\\
F_W &= \frac{c_W\,P^*_h}{W\,(W^2-m_0^2)}\ \Bigl[+\ \frac{a_R\,M_R^2\Gamma_R^2}{(W^2-M_R^2)^2+M_R^2\Gamma_R^2}\ \ (\Lambda)\Bigr],\qquad
\sigma_\mathrm{fact} = F_Q\,F_t\,F_W,
\lbl{eq:sigfact}
\end{feqa}
with the parameters of \tabr{tab:sigfact}. The $\Lambda$ branch is taken for $M_Y<1150$ MeV, the $\Sigma$ branch
otherwise. The fit is a function of $W$, $Q^2$ and $t$ only: there is no $\varepsilon$ or $\varphi$ dependence, so the
longitudinal--transverse separation and the interference terms are not modelled. The proton mass enters $\nu_W$ for
every final state (also for $\Sigma^-$ on the neutron and for coherent production), $m_0$ is a hard-coded proton mass
in GeV, and for coherent production $W$ is the photon--nucleus mass (4--13.5 GeV), far outside the range of a nucleon
fit. As a check, $P^*_h$ is recomputed from $W$ and $M_Y$; a message is printed when the two differ by more than
0.1\,\% (this happens for the central event of the $\Sigma$ and nuclear runs). If $W^2=0$ (``central event'' in the
code comment) $F_W$ is not computed and keeps its previous value. The coefficients are transcribed from the code and
were not compared with the thesis.

\begin{simctable}{footnotesize}{L{0.14\linewidth}L{0.2\linewidth}L{0.2\linewidth}L{0.46\linewidth}}
\caption{Parameters of the factorized model \fsrc{physics\_kaon.f:215-232} (one parameter per row; GeV units).}\lbl{tab:sigfact}\\
\toprule\rowcolor{simchead}\simcth{parameter} & \simcth{$\Lambda$ ($M_Y<1150$ MeV)} & \simcth{$\Sigma^0$, $\Sigma^-$, coherent} & \simcth{meaning}\\\midrule\endfirsthead
\toprule\rowcolor{simchead}\simcth{parameter} & \simcth{$\Lambda$ ($M_Y<1150$ MeV)} & \simcth{$\Sigma^0$, $\Sigma^-$, coherent} & \simcth{meaning}\\\midrule\endhead
\bottomrule\endlastfoot
$A_Q$ & 2.67 GeV$^2$ & 0.79 GeV$^2$ & dipole-like $Q^2$ scale\\
$b$ & 2.1 GeV$^{-2}$ & 1.0 GeV$^{-2}$ & $t$ slope\\
$c_W$ & $0.959\times4.1959$ & $0.959\times4.1959$ & normalisation of the $W$ factor\\
$m_0$ & 0.93827 GeV & 0.93827 GeV & proton mass in the phase-space denominator\\
$a_R$ & 0.18 & -- (term commented out) & resonance strength\\
$M_R$ & 1.72 GeV & -- & resonance mass\\
$\Gamma_R$ & 0.10 GeV & -- & resonance width\\
\end{simctable}

\paragraph{Saghai amplitude model (diagnostic column \code{saghai}).} For each event the model also evaluates a
cross section from tabulated CGLN-type amplitudes $F_1\ldots F_6$ (complex) attributed in the code to Saghai
(``Saghai's form'', \fsrc{physics\_kaon.f:327}); the code cites no paper. The amplitudes are interpolated in
$(s, Q^2, \theta^*)$ from \code{saghai\_proton.dat} ($\Lambda$) or \code{saghai\_sigma0.dat} ($\Sigma$), read at start-up
(\tabr{tab:saghai}), with $s=W^2_\mathrm{cm}$ in GeV$^2$, $Q^2$ in GeV$^2$ and $\theta^*$ in degrees. The interpolation is
the CERNLIB multilinear \code{fint}:
\begin{feqa}{physics/kaon/KaonCrossSections.cpp: fint}{cern/fint.f:27-70}
F(x_1,x_2,x_3) &= \sum_{\epsilon\in\{0,1\}^3}\ \prod_{n=1}^{3}\bigl[(1-\eta_n)\,\delta_{\epsilon_n0}+\eta_n\,\delta_{\epsilon_n1}\bigr]\,
T(k_1+\epsilon_1,k_2+\epsilon_2,k_3+\epsilon_3),\nonumber\\
\eta_n &= \frac{x_n-g_{n,k_n}}{g_{n,k_n+1}-g_{n,k_n}},
\lbl{eq:fint}
\end{feqa}
$k_n$ being the grid interval of $x_n$ found by bisection and limited to the first and last interval (linear
extrapolation outside the grid); on an exact node the dimension is not interpolated. With
$x=\cos\theta^*$, the kinematic factors and the four response terms are
\begin{feqa}{physics/kaon/KaonCrossSections.cpp: eekeek, eekeeks}{physics\_kaon.f:264-280, 320-348; 377-393, 452-474}
P_K^* &= \frac{\sqrt{\max\bigl(0,(W^2-M_K^2-M_Y^2)^2-4M_K^2M_Y^2\bigr)}}{2W},\qquad
\nu^*_\gamma = \frac{M_p\nu_W-Q^2}{W},\qquad r = \frac{Q}{\nu^*_\gamma},\nonumber\\
a &= \frac{P^*_K}{2W\,(W^2-M_p^2)}\,(\hbar c)^2\cdot10^4,\qquad a_L = a\,r^2,\nonumber\\
\sigma_T &= a\Bigl[|F_1|^2+|F_2|^2+2x\,\mathrm{Re}(F_1^*F_2)+\tfrac12\sin^2\theta^*\bigl(|F_3|^2+|F_4|^2+2\,\mathrm{Re}(F_1^*F_4-F_2^*F_3+x\,F_3^*F_4)\bigr)\Bigr],\nonumber\\
\sigma_L &= a_L\,\varepsilon\,\bigl[|F_5|^2+|F_6|^2+2x\,\mathrm{Re}(F_5^*F_6)\bigr],\nonumber\\
\sigma_{TT} &= a\,\varepsilon\,\sin^2\theta_{pq}\cos2\varphi\,\bigl[\tfrac12|F_3|^2+\tfrac12|F_4|^2+\mathrm{Re}(F_1^*F_4-F_2^*F_3+x\,F_3^*F_4)\bigr],\nonumber\\
\sigma_{LT} &= a\sqrt{2r^2\varepsilon(1+\varepsilon)}\,\sin\theta_{pq}\cos\varphi\ \mathrm{Re}\bigl[F_5(F_3^*-F_2^*+xF_4^*)+F_6(F_1^*+xF_3^*+F_4^*)\bigr],\nonumber\\
\sigma_\mathrm{Saghai} &= \sigma_T+\sigma_L+\sigma_{TT}+\sigma_{LT},
\lbl{eq:saghai}
\end{feqa}
($\nu_W$ as in \eqr{eq:sigfact}; the factor $(\hbar c)^2\cdot10^4$ converts fm$^2$ to $\mu$b, so $\sigma_\mathrm{Saghai}$ is in $\mu$b/sr if the tabulated amplitudes are in fm, which the code does not state; the $\Sigma^0$ routine is identical except for the table and its grid). Three remarks on
the code as written: the interference terms use the \emph{laboratory} angle $\theta_{pq}$ where $\sigma_T$ uses
$\theta^*$; the azimuth $\varphi$ passed to the routine is never assigned and is 0 (the code comment says ``we are always
calculating for $\varphi=0$''); and compared with the CGLN response functions in the convention of Drechsel and Tiator
\citer{dt}, the expressions correspond to $F_2\to-F_2$ (and a sign convention of the LT term), which may be the convention
of the tables; we did not check Saghai's convention. The routine is called with $\Lambda$ tables for $M_Y<1150$ MeV
and with the $\Sigma$ tables otherwise, so coherent production (large $M_Y$) also uses the $\Sigma$ tables.

\begin{simctable}{footnotesize}{L{0.2\linewidth}L{0.17\linewidth}L{0.63\linewidth}}
\caption{The Saghai tables and their grids \fsrc{physics\_kaon.f:281-302, 394-435; dbase.f:650-683} (one grid per row).}\lbl{tab:saghai}\\
\toprule\rowcolor{simchead}\simcth{table} & \simcth{variable} & \simcth{nodes}\\\midrule\endfirsthead
\toprule\rowcolor{simchead}\simcth{table} & \simcth{variable} & \simcth{nodes}\\\midrule\endhead
\bottomrule\endlastfoot
\code{saghai\_proton.dat} & $s$ [GeV$^2$] & 10 nodes $2.6+0.3(i-1)$ in the compat mode; the file holds blocks at 5.1 and 5.4 for nodes 9 and 10, which the default mode uses\\
\code{saghai\_proton.dat} & $Q^2$ [GeV$^2$] & 11 nodes, 0 to 2.0 in steps of 0.2\\
\code{saghai\_proton.dat} & $\theta^*$ [deg] & 19 nodes, 0 to 180 in steps of 10\\
\code{saghai\_sigma0.dat} & $s$ [GeV$^2$] & 20 nodes: 2.851, 2.898, 2.945, 2.991, 3.038, 3.085, 3.132, then 3.320 to 5.572 in steps of about 0.188\\
\code{saghai\_sigma0.dat} & $Q^2$ [GeV$^2$] & 10 nodes: 0, 0.25, 0.376, 0.52, 0.75, 1.0, 1.25, 1.5, 1.75, 2.0\\
\code{saghai\_sigma0.dat} & $\theta^*$ [deg] & 19 nodes, 0 to 180 in steps of 10\\
\code{weightc.dat} & -- & $20\times50$ phase-space weights for HMS/SOS (read, used only by an uncalled routine)\\
\code{weightd.dat} & -- & $8\times40\times30$ phase-space weights for SOS/HMS (read, used only by an uncalled routine)\\
\end{simctable}
Each proton record has a kinematics line and two lines of three complex amplitudes; the $\Sigma^0$ file starts with a
truncated line, so a plain sequential read takes every record one line late. The compat mode reproduces that read and
also the interpolation result as the Fortran build returns it, which is $+0$ for every value (the tree column
\code{saghai} is then 0); the default mode skips the stray line, uses the grid of the file and returns the interpolated
value. None of this enters the weight.

\subsection{Weight and units}
\begin{feq}{reactions/kaon/KaonProduction.cpp: KaonProduction::completeMain}{event.f:1505-1512, 1536-1538, 1561-1564}
w = J\cdot w_\mathrm{gen}\cdot\sigma\cdot n_\mathrm{tgt}\cdot\bigl(1+V_C/E_b\bigr)^2_{[Z>1]}\ \bigl[\cdot P_\mathrm{surv}\ \text{if the decay is not simulated}\bigr],
\lbl{eq:weight}
\end{feq}
with $n_\mathrm{tgt}$ of \tabr{tab:masses} ($Z$ or $N$ also for coherent production) and $P_\mathrm{surv}$ of
\eqr{eq:surv}. $w$ is in $\mu$b/(MeV\,sr$^2$), which matches $V_\mathrm{gen}$ in sr$^2$ MeV.

\paragraph{Central event.}\lbl{sec:central} The \texttt{.hist} line \code{CENTRAL.sigcc} evaluates the model once for the
central momenta and angles of both arms. In the default mode \code{prepareCentral} first gives the central event its
$W$, $\varepsilon$, $\theta_{pq}$, $-t$, $-t_\mathrm{min}$ and $\varphi_{pq}$ (\eqr{eq:mesonvars}); in the compat mode
these stay 0, as in the Fortran, where the model then sees $\varepsilon=0$ and $t=0$. The focal-plane path of the
central event is 0, so its survival probability is 1. The central kaon is placed at $P_h^0$ whatever the hyperon, which
is why the $P^*_h$ check of \eqr{eq:sigfact} prints its message for $\Sigma$ and nuclear runs.

% =====================================================================================================================
\section{Radiation, energy loss and decays}\lbl{sec:rad}
\subsection{Radiation}
All three tails radiate; the kaon is the radiating hadron, with $\lambda_3$ from its own $E_h$, $P_h$
(\ceqr{eq:lambda}). The photon limits are those of the ``meson'' case of \code{generateRad}:
\begin{feqa}{radiation/Radc.cpp: generateRad}{radc.f:297-308, 372-389, 428-446}
E_{\gamma1} &\in \bigl[-d_E,\ \min(\Sigma_\mathrm{max}-E',\ E^\mathrm{max}_{\gamma1})+d_E\bigr],\nonumber\\
E_{\gamma2} &\in \bigl[E'-E'_\mathrm{max}-d_E,\ \min(E'-E'_\mathrm{min},\ E_\gamma^\mathrm{tot,max}-E_{\gamma1})+d_E\bigr],\nonumber\\
E_{\gamma3} &\in \bigl[E_h-E_{h,\mathrm{max}}-d_E,\ \min(E_h-E_{h,\mathrm{min}},\ E_\gamma^\mathrm{tot,max}-E_{\gamma1}-E_{\gamma2})+d_E\bigr].
\lbl{eq:radlimits}
\end{feqa}
After a tail-1 photon \code{completeEv} runs again at $E_\mathrm{in}-E_{\gamma1}$ with the same directions and the same
Fermi momentum: a new kaon energy (\eqr{eq:quad}) and new meson variables; the cross section is evaluated there. The
soft and hard factors use the elastic e--p kinematics of \ceqr{eq:bremkin}, which are not those of kaon production.

\subsection{Energy loss and multiple scattering}
The kaon loses energy and scatters with its own mass and $\beta=P_h/E_h$ like any hadron (\csecr{sec:eloss}); the
reconstruction adds back the most probable loss with \code{correct\_Eloss}.

\subsection{Kaon decay in flight (\code{doing\_decay} = 1)}
The spectrometer models check the kaon for decay in every drift and every transport element, with the input $c\tau$
(371.2 cm for the K$^+$ in the reference inputs). For a step of nominal length $L$ (the transport element adds its
path-length correction $\Delta L$) and the current momentum $P$:
\begin{feqa}{optics/CosyEngine.cpp: project, decayProducts}{shared/transp.f:134-185, 231-276; shared/project.f:43-117}
\beta\gamma &= P/M,\qquad z_\mathrm{dec} = -c\tau\,\beta\gamma\,\ln(1-\xi_1),\nonumber\\
&\text{decay in the step if } z_\mathrm{dec}\le L\sqrt{1+x'^2+y'^2}\ \text{(drift)},\ \ z_\mathrm{dec}\le L+\Delta L\ \text{(element)},\nonumber\\
\cos\theta^\star &= 2\xi_2-1,\quad \varphi^\star = 2\pi\xi_3,\qquad
(p^\star,\ m_f) = \begin{cases}(235.5\ \mathrm{MeV}/c,\ m_\mu) & \xi_4<0.7\ \ (K\to\mu\nu)\\
\bigl(\sqrt{M_K^2/4-M_\pi^2},\ M_\pi\bigr) & \text{otherwise}\ \ (K\to\pi\pi)\end{cases}\nonumber\\
(E,\vec p) &= \Lambda(-\beta\hat t)\,\bigl(\sqrt{m_f^2+p^{\star2}},\ p^\star\hat n^\star\bigr),\qquad
x' = p_x/p_z,\quad y'=p_y/p_z,\quad \delta = 100\,(|\vec p|/P^0-1),
\lbl{eq:decay}
\end{feqa}
$\hat t$ being the track direction and $\hat n^\star$ the isotropic direction in the kaon rest frame. A transport element
applies a decay found in its first half before the transport of the element and one in its second half after it. The
decay product continues with its own mass (used for multiple scattering and any later energy loss), reported in the tree
column \code{Mhadron}; it does not decay again (a pion from $K\to\pi\pi$ is not re-tested for its own decay after a kaon
decay). The tree column \code{decdist} is the distance travelled before the decay (the whole path if the kaon did not
decay). The model is fixed: 70\,\% $\mu\nu$ and 30\,\% $\pi\pi$ (the measured K$^+$ branching ratios are 63.6\,\% and
20.7\,\% for $\pi^+\pi^0$, with further three-body modes), and the $\pi\pi$ momentum uses the equal-mass formula with the
charged pion, 203.6 MeV/$c$ (205.1 MeV/$c$ for $\pi^+\pi^0$). In a drift, the default mode ends a decayed track at the
end of the drift; the compat mode reproduces the Fortran, which computes the remaining drift with the slopes after the
decay and so ends the track slightly off.

Up to SIMC 294f662 the $K\to\pi\pi$ product of a decay in the first half of a transport element kept the kaon mass;
upstream fixed that line (the daughter mass is $M_\pi$ since 59ccdf6), and SIMC++ uses $M_\pi$ in both modes. The
option \code{-{}-compat-target 294f662} restores the kaon mass for the old references.

\subsection{Survival probability (\code{doing\_decay} = 0)}
Without the decay simulation the weight is multiplied by the probability that the kaon reaches the detector that
identifies it:
\begin{feqa}{physics/kaon/KaonCrossSections.cpp: peeK}{physics\_kaon.f:148-164}
P_\mathrm{surv} &= \exp\Bigl[-\frac{L}{c\tau\,\beta\gamma}\Bigr],\qquad
L = L_\mathrm{FP} + z_\mathrm{aero}\sqrt{1+x'^2_\mathrm{fp}+y'^2_\mathrm{fp}}\ \ \bigl[\text{compat: }z_\mathrm{aero}(1+x'^2_\mathrm{fp}+y'^2_\mathrm{fp})\bigr],\nonumber\\
\beta\gamma &= P/M_K,\qquad P = P_h^0\,(1+\delta^\mathrm{SP}_h/100)\ \ \bigl[\text{compat: } P=P_h^0\bigr],
\lbl{eq:surv}
\end{feqa}
where $L_\mathrm{FP}$ is the path length of the hadron-arm model (to the back of its last detector), $x'_\mathrm{fp}$,
$y'_\mathrm{fp}$ the focal-plane slopes, and $z_\mathrm{aero}$ the hard-coded distance of the aerogel detector from the
end of the path: 0 for the HMS, $-82.8$ cm for the SOS, $-183$ cm for the HRS; for the SHMS the compat mode keeps the
value of the previous call (0 in a run) and the default mode uses 0. $\delta^\mathrm{SP}_h$ is the kaon at the
spectrometer entrance (after radiation and energy loss): the code comment asks for that momentum, but the Fortran uses
the central momentum, which makes $P_\mathrm{surv}$ almost constant. The value is written to \code{decdist}.

% =====================================================================================================================
\section{Inputs and outputs}\lbl{sec:io}
\subsection{Inputs}
\begin{simctable}{footnotesize}{L{0.24\linewidth}L{0.12\linewidth}L{0.08\linewidth}L{0.56\linewidth}}
\caption{Inputs specific to kaon production (one input per row; the common inputs are those of note 01, \S7, and the
user manual, chapter 4).}\lbl{tab:inputs}\\
\toprule\rowcolor{simchead}\simcth{input} & \simcth{unit} & \simcth{default} & \simcth{meaning here}\\\midrule\endfirsthead
\toprule\rowcolor{simchead}\simcth{input} & \simcth{unit} & \simcth{default} & \simcth{meaning here}\\\midrule\endhead
\bottomrule\endlastfoot
\code{doing\_kaon} & flag & 0 & 1 selects this reaction\\
\code{which\_kaon} & code & 0 & final state (\tabr{tab:masses})\\
\code{targ\%A} & -- & 0 & 1: hydrogen; 2: deuterium; $\ge3$: nucleus (momentum distribution, spectral function)\\
\code{targ\%Z} & -- & 0 & $n_\mathrm{tgt}$ for $\Lambda$, $\Sigma^0$; Coulomb correction for $Z>1$\\
\code{targ\%mass\_amu} & u & 0 & $M_A$ (forced to $M_p$, $M_d$ for $A=1,2$)\\
\code{targ\%mrec\_amu} & u & 0 & $M_\mathrm{rec}$ of the spectator system ($A\ge3$); ignored for coherent production\\
\code{doing\_decay} & flag & 0 & 1: decay in flight (\eqr{eq:decay}); 0: survival weight (\eqr{eq:surv})\\
\code{ctau} & cm & 0 & $c\tau$ of the kaon (371.2 cm for K$^+$); used in both cases\\
\code{hadron\_arm} & code & 0 & the kaon arm; also selects $z_\mathrm{aero}$\\
\code{spec\%p\%P} & MeV/$c$ & 0 & $P_h^0$; the survival $\beta\gamma$ in the compat mode\\
\code{doing\_semi} & flag & 0 & with \code{doing\_kaon} = 1 the run is kaon SIDIS instead (another note)\\
\code{doing\_phsp} & flag & 0 & 1: phase-space run with the kaon mass\\
\end{simctable}

\subsection{Outputs}
The sink \code{root\_tree} writes one entry of \code{h10} per success with the 66 columns of \tabr{tab:columns} (layout
of upstream SIMC 59ccdf6; \code{-{}-tree-layout legacy} gives the 55 columns of SIMC 294f662). The upstream Fortran
writes 66 values but declares only 65 column names for kaon runs and stops at the first event; SIMC++ declares the
column \code{sigcm} as well, in both modes. Columns 1--24 are arranged as in the (e,e\pr p) layout: \code{hs} is the
arm on the beam right, \code{ss} the arm on the beam left.

\begin{simctable}{footnotesize}{L{0.06\linewidth}L{0.15\linewidth}L{0.13\linewidth}L{0.66\linewidth}}
\caption{Tree columns of kaon production (one column per row).}\lbl{tab:columns}\\
\toprule\rowcolor{simchead}\simcth{\#} & \simcth{name} & \simcth{unit} & \simcth{content}\\\midrule\endfirsthead
\toprule\rowcolor{simchead}\simcth{\#} & \simcth{name} & \simcth{unit} & \simcth{content}\\\midrule\endhead
\bottomrule\endlastfoot
1 & \code{hsdelta} & \% & reconstructed $\delta$, \code{hs} arm\\
2 & \code{hsyptar} & rad & reconstructed $y'$, \code{hs} arm\\
3 & \code{hsxptar} & rad & reconstructed $x'$, \code{hs} arm\\
4 & \code{hsytar} & cm & reconstructed horizontal target position, \code{hs} arm\\
5 & \code{hsxfp} & cm & focal-plane $x$, \code{hs} arm\\
6 & \code{hsxpfp} & rad & focal-plane $x'$, \code{hs} arm\\
7 & \code{hsyfp} & cm & focal-plane $y$, \code{hs} arm\\
8 & \code{hsypfp} & rad & focal-plane $y'$, \code{hs} arm\\
9 & \code{hsdeltai} & \% & true $\delta$ after radiation, \code{hs} arm\\
10 & \code{hsyptari} & rad & true $y'$, \code{hs} arm\\
11 & \code{hsxptari} & rad & true $x'$, \code{hs} arm\\
12 & \code{hsytari} & cm & $z_t\sin\theta^0$ of the \code{hs} arm\\
13 & \code{ssdelta} & \% & reconstructed $\delta$, \code{ss} arm\\
14 & \code{ssyptar} & rad & reconstructed $y'$, \code{ss} arm\\
15 & \code{ssxptar} & rad & reconstructed $x'$, \code{ss} arm\\
16 & \code{ssytar} & cm & reconstructed horizontal target position, \code{ss} arm\\
17 & \code{ssxfp} & cm & focal-plane $x$, \code{ss} arm\\
18 & \code{ssxpfp} & rad & focal-plane $x'$, \code{ss} arm\\
19 & \code{ssyfp} & cm & focal-plane $y$, \code{ss} arm\\
20 & \code{ssypfp} & rad & focal-plane $y'$, \code{ss} arm\\
21 & \code{ssdeltai} & \% & true $\delta$ after radiation, \code{ss} arm\\
22 & \code{ssyptari} & rad & true $y'$, \code{ss} arm\\
23 & \code{ssxptari} & rad & true $x'$, \code{ss} arm\\
24 & \code{ssytari} & cm & $-z_t\sin\theta^0$ of the \code{ss} arm\\
25 & \code{q} & GeV/$c$ & $|\vec q^\mathrm{rec}|$\\
26 & \code{nu} & GeV & $\nu^\mathrm{rec}$\\
27 & \code{Q2} & (GeV/$c$)$^2$ & $(Q^\mathrm{rec})^2$\\
28 & \code{W} & GeV & $W^\mathrm{rec}$ (proton mass)\\
29 & \code{xbj} & -- & $x^\mathrm{rec}$\\
30 & \code{epsilon} & -- & $\varepsilon^\mathrm{rec}$\\
31 & \code{Em} & GeV & $E_m^\mathrm{rec}$ of \eqr{eq:recon}\\
32 & \code{Pm} & GeV/$c$ & $P_m^\mathrm{rec}$\\
33 & \code{qi} & GeV/$c$ & vertex $|\vec q|$\\
34 & \code{nui} & GeV & vertex $\nu$\\
35 & \code{Q2i} & (GeV/$c$)$^2$ & vertex $Q^2$\\
36 & \code{Wi} & GeV & vertex $W$ of \eqr{eq:mesonvars} (without the Fermi momentum)\\
37 & \code{xbji} & -- & vertex $x$\\
38 & \code{epsiloni} & -- & vertex $\varepsilon$\\
39 & \code{Emi} & GeV & vertex $E_\mathrm{miss}=\nu+M_A-E_h$\\
40 & \code{Pmi} & GeV/$c$ & vertex $P_\mathrm{miss}$\\
41 & \code{missmass} & GeV & $M_x$ of \eqr{eq:recon}\\
42 & \code{mmnuc} & GeV & $M_x^A$ of \eqr{eq:recon}\\
43 & \code{phad} & GeV/$c$ & $P_h^\mathrm{rec}$\\
44 & \code{t} & GeV$^2$ & $-t^\mathrm{rec}$\\
45 & \code{ti} & GeV$^2$ & vertex $-t$\\
46 & \code{z} & -- & $E_h^\mathrm{rec}/\nu^\mathrm{rec}$\\
47 & \code{zi} & -- & vertex $E_h/\nu$\\
48 & \code{pt2} & GeV$^2$ & reconstructed $p_T^2$\\
49 & \code{pt2i} & GeV$^2$ & vertex $p_T^2$\\
50 & \code{thetapq} & rad & $\theta^\mathrm{rec}_{pq}$\\
51 & \code{phipq} & rad & $\varphi^\mathrm{rec}_{pq}$\\
52 & \code{thetapqi} & rad & vertex $\theta_{pq}$\\
53 & \code{phipqi} & rad & vertex $\varphi_{pq}$\\
54 & \code{pmpar} & GeV/$c$ & $P^\mathrm{rec}_{m,\parallel}$\\
55 & \code{pmper} & GeV/$c$ & $P^\mathrm{rec}_{m,\perp}$\\
56 & \code{pmoop} & GeV/$c$ & $P^\mathrm{rec}_{m,\mathrm{oop}}$\\
57 & \code{fry} & cm & minus the vertical raster coordinate\\
58 & \code{radphot} & GeV & $E_{\gamma1}+E_{\gamma2}+E_{\gamma3}$\\
59 & \code{pfermi} & GeV/$c$ & $p_F$ with the sign of $\hat p_F\cdot\hat q$ (0 without Fermi motion)\\
60 & \code{siglab} & $\mu$b/MeV/sr$^2$ & $\sigma$ of \eqr{eq:peek} with the Coulomb factor\\
61 & \code{sigcm} & $\mu$b/sr & $\sigma_\mathrm{fact}$ (the model in the weight)\\
62 & \code{saghai} & see text & $\sigma_\mathrm{Saghai}$ of \eqr{eq:saghai} (0 in the compat mode)\\
63 & \code{factor} & $\mu$b/sr & $\sigma_\mathrm{fact}$ again (the factorized model by name)\\
64 & \code{Weight} & $\mu$b/MeV/sr$^2$ & $w$ of \eqr{eq:weight}\\
65 & \code{decdist} & cm or -- & decay distance (decay on) or $P_\mathrm{surv}$ (decay off)\\
66 & \code{Mhadron} & MeV & mass of the particle that reached the detectors ($M_K$, $m_\mu$ or $M_\pi$)\\
\end{simctable}

% =====================================================================================================================
\section{Validation in SIMC++}\lbl{sec:validation}
\begin{simctable}{footnotesize}{L{0.5\linewidth}L{0.22\linewidth}L{0.28\linewidth}}
\caption{CTest tests that cover kaon production (one test per row).}\lbl{tab:tests}\\
\toprule\rowcolor{simchead}\simcth{test} & \simcth{labels} & \simcth{what it checks}\\\midrule\endfirsthead
\toprule\rowcolor{simchead}\simcth{test} & \simcth{labels} & \simcth{what it checks}\\\midrule\endhead
\bottomrule\endlastfoot
\code{l1\_kaon\_data\_files} & l1 & the four tables are the Fortran files\\
\code{l1\_saghai\_tables} & l1 & the tables as read (both readings of the $\Sigma^0$ file)\\
\code{l1\_sig\_factorized} & l1 & \eqr{eq:sigfact} against the Fortran table\\
\code{l1\_eekeek} & l1 & \eqr{eq:saghai}, $\Lambda$, default numerics\\
\code{l1\_eekeek\_fintabi} & l1 & the same with the Fortran interpolation result\\
\code{l1\_eekeeks} & l1 & \eqr{eq:saghai}, $\Sigma$\\
\code{l1\_eekeeks\_fintabi} & l1 & the same with the Fortran interpolation result\\
\code{l1\_phspwght} & l1 & the uncalled phase-space weight routine\\
\code{l1\_phspwght\_fintabi} & l1 & the same, Fortran interpolation result\\
\code{l1\_peek} & l1 & \eqr{eq:peek} and \eqr{eq:surv}\\
\code{l1\_peek\_fintabi} & l1 & the same, Fortran interpolation result\\
\code{l1\_kaon\_plugin} & l1 & the model through its plug-in interface\\
\code{l1\_transform\_to\_cm\_kaon} & l1 & \eqsr{eq:cm}{eq:jold}\\
\code{l1\_kaon\_state} & l1 & set-up values (masses, limits)\\
\code{l1\_generate\_kaon\_fixedrng} & l1 & \eqsr{eq:genslopes}{eq:fermi} with fixed random numbers\\
\code{l1\_complete\_ev\_kaon} & l1 & \eqsr{eq:quad}{eq:trec}\\
\code{l1\_complete\_recon\_ev\_kaon} & l1 & \eqr{eq:recon}\\
\code{l1\_complete\_main\_kaon} & l1 & \eqr{eq:weight}\\
\code{l1\_decay\_kaon\_fixedrng} & l1 & \eqr{eq:decay} in the drifts\\
\code{l1\_mc\_decay\_kaon\_fixedrng} & l1 & decays inside the spectrometer models\\
\code{l1\_replay\_event\_capture\_kaon\_kl\_h} & l1, replay & every routine of captured H $\Lambda$ events\\
\code{l1\_replay\_event\_capture\_kaon\_ksm\_d} & l1, replay & every routine of captured D $\Sigma^-$ events\\
\code{l1\_simcrun\_replay\_kaon\_kl\_h} & l1, replay & complete H run from a captured random stream, bitwise\\
\code{l1\_simcrun\_replay\_kaon\_ksm\_d} & l1, replay & complete D run, bitwise\\
\code{l1\_montecarlo\_replay\_kaon\_kl\_h} & l1, replay & the spectrometer stage of the H capture\\
\code{l1\_montecarlo\_replay\_kaon\_ksm\_d} & l1, replay & the spectrometer stage of the D capture\\
\code{l1\_simcrun\_ranlux\_kaon\_h\_k1\_norad\_dec0\_A} & l1, replay, slow & full run with RANLUX, compat mode, against the 294f662 Fortran run\\
\code{l1\_simcrun\_ranlux\_up\_upf1\_kaon\_h\_k1\_rad\_dec1\_nontu} & l1, replay, upstream\_layout, slow & decay-on run against the 59ccdf6 Fortran (no tree: the Fortran stops when writing it)\\
\code{regression\_identity\_kaon\_h\_k1\_rad\_dec1} & regression & compat-mode outputs against the baseline\\
\code{regression\_corrected\_kaon\_h\_k1\_rad\_dec1} & regression & default-mode outputs against the baseline\\
\code{regression\_identity\_kaon\_d\_k1\_sigm\_rad\_dec1} & regression & compat-mode $\Sigma^-$ D outputs\\
\code{regression\_corrected\_kaon\_d\_k1\_sigm\_rad\_dec1} & regression & default-mode $\Sigma^-$ D outputs\\
\code{regression\_upstream\_u59\_kaon\_d\_k1\_sigm\_rad\_dec1\_A\_nontu} & regression\_upstream, slow & decay-on D run against 59ccdf6\\
\code{regression\_upstream\_u59\_shipped\_test\_kaon\_d\_pct} & regression\_upstream, slow & a kaon run with the tree, against the Fortran run without it\\
\code{fortran\_diff\_D043\_kaon\_model} & fortran\_diff, unit & the model corrections in both modes\\
\code{fortran\_diff\_D043\_F13\_central\_sigcc\_kaon} & fortran\_diff, unit & the central event in both modes\\
\code{fortran\_diff\_D006\_kaon\_decay\_mass} & fortran\_diff, unit & the $K\to\pi\pi$ daughter mass for both compat targets\\
\code{fortran\_diff\_D037\_F7\_project\_decay\_drift} & fortran\_diff, unit & the drift after a decay in both modes\\
\code{fortran\_diff\_U\_A2\_A4\_upstream\_layout\_kaon\_pi0\_field} & fortran\_diff, unit & the 66 kaon columns\\
\end{simctable}
The validation ledger (\code{project\_control/VALIDATION\_LEDGER.md}, entry of 2026-10-11) records all kaon L1 and replay
tests bitwise and the statistical comparison of the compat mode with the nine Fortran reference runs passed
(\code{validation/simc\_compare}, seed 4242, $2\times10^5$ events; one seed-4242 fluctuation, other seeds pass); a
RANLUX compat run reproduces \code{kaon\_h\_k1\_norad\_dec0\_A} exactly (.hist, .geni, .gen and the tree of 200000
entries $\times$ 55 columns). FORTRAN\_DIFFERENCES.md \S10 lists, per run, the effect of each default-mode correction.

% =====================================================================================================================
\section{Fortran origin and Fortran findings}\lbl{sec:origin}
\begin{simctable}{footnotesize}{L{0.42\linewidth}L{0.58\linewidth}}
\caption{Fortran routine $\to$ C++ function (one routine per row; the common routines are in \ctabr{tab:map}).}\lbl{tab:map}\\
\toprule\rowcolor{simchead}\simcth{Fortran (upstream 59ccdf6)} & \simcth{C++ (src/...)}\\\midrule\endfirsthead
\toprule\rowcolor{simchead}\simcth{Fortran (upstream 59ccdf6)} & \simcth{C++ (src/...)}\\\midrule\endhead
\bottomrule\endlastfoot
\file{dbase.f:163}-180 (kaon flags) & \csrc{core/RunSetup.cpp: DeriveReactionFlags}\\
\file{dbase.f:400}-442 (kaon masses) & \csrc{setup/SetupReaction.cpp: BuildReactionAndKinematics}\\
\file{dbase.f:567}-591 (momentum distribution) & \csrc{setup/SetupTheory.cpp: ReadMomentumDistribution}\\
\file{dbase.f:598}-628 (spectral function for $E_m$) & \csrc{physics/nuclear/SpectralFunctionTables.cpp: SfTable::load}\\
\file{dbase.f:630}-685 (kaon tables) & \csrc{physics/kaon/KaonCrossSections.cpp: SaghaiTables::load}\\
\file{init.f:344}-510 (limits) & \csrc{setup/SetupLimits.cpp: BuildLimits}\\
\file{event.f} generate, l.276-372 & \csrc{reactions/pion/PionProduction.cpp: PionProduction::generate}\\
\file{sf\_lookup.f} generate\_em & \csrc{physics/nuclear/SpectralFunctionTables.cpp: SfTable::generateEm}\\
\file{event.f} complete\_ev, l.608-698, 947-951, 1015-1025 & \csrc{reactions/pion/PionProduction.cpp: PionProduction::completeEv}\\
\file{event.f} complete\_ev, l.723-771 & \csrc{reactions/pion/PionProduction.cpp: PionProduction::pionMainKinematics}\\
\file{event.f} complete\_recon\_ev, l.1325-1354 & \csrc{reactions/EepReactionBase.cpp: EepReactionBase::completeRecon}\\
\file{event.f} complete\_main, l.1505-1512, 1561-1564 & \csrc{reactions/kaon/KaonProduction.cpp: KaonProduction::completeMain}\\
\file{jacobians.f} transform\_to\_cm & \csrc{reactions/EventChain.cpp: transformToCm}\\
\file{physics\_kaon.f} peeK & \csrc{physics/kaon/KaonCrossSections.cpp: peeK}\\
\file{physics\_kaon.f} sig\_factorized & \csrc{physics/kaon/KaonCrossSections.cpp: sigFactorized}\\
\file{physics\_kaon.f} eekeek & \csrc{physics/kaon/KaonCrossSections.cpp: eekeek}\\
\file{physics\_kaon.f} eekeeks & \csrc{physics/kaon/KaonCrossSections.cpp: eekeeks}\\
\file{physics\_kaon.f} phspwght (not called) & \csrc{physics/kaon/KaonCrossSections.cpp: phspwght}\\
\file{cern/fint.f} & \csrc{physics/kaon/KaonCrossSections.cpp: fint}\\
\file{physics\_kaon.f} (call of peeK) & \csrc{physics/kaon/KaonModel.cpp: KaonModel::sigccMain}\\
\file{simc.f} calculate\_central, l.1153-1316 & \csrc{reactions/pion/PionProduction.cpp: PionProduction::prepareCentral}\\
\file{shared/transp.f:134}-276 (decay) & \csrc{optics/CosyEngine.cpp: decayProducts}\\
\file{shared/project.f} & \csrc{optics/CosyEngine.cpp: project}\\
\file{radc.f:297}-308 & \csrc{radiation/Radc.cpp: generateRad}\\
\file{results\_write.f:178}-234 & \csrc{output/NtupleLayout.cpp: UpstreamMesonColumns}\\
\file{simc.f:383}-397 & \csrc{reactions/EepReactionBase.cpp: EepReactionBase::genVolume}\\
\end{simctable}

\begin{simcfindings}{footnotesize}{L{0.2\linewidth}L{0.48\linewidth}L{0.32\linewidth}}
\caption{Fortran findings that touch kaon production (one finding per row; the common ones are in \ctabr{tab:findings}).}\lbl{tab:findings}\\
\toprule\rowcolor{simchead}\simcth{id} & \simcth{Fortran behaviour} & \simcth{default / compat}\\\midrule\endfirsthead
\toprule\rowcolor{simchead}\simcth{id} & \simcth{Fortran behaviour} & \simcth{default / compat}\\\midrule\endhead
\bottomrule\endlastfoot
\finding{D-043 F1} & the single-precision callers of \code{fint} read 0 for every value: \code{saghai} = 0 & interpolated value / 0\\
\finding{D-043 F2} & \code{saghai\_sigma0.dat} read one line off & records read in place / Fortran\\
\finding{D-043 F3} & nodes 9, 10 of the $\Lambda$ $s$ grid are 5.0, 5.3 but the file blocks are 5.1, 5.4 & file grid / Fortran grid\\
\finding{D-043 F5} & survival $\beta\gamma$ from the central momentum & kaon momentum / central\\
\finding{D-043 F6} & SHMS $z_\mathrm{aero}$ stale; path factor without square root & 0 and square root / Fortran\\
\finding{D-043 F13} & the central event sees $\varepsilon=t=\theta_{pq}=\varphi_{pq}=0$ & central kinematics / zeros\\
\finding{D-006 kaon decay} & $K\to\pi\pi$ daughter with the kaon mass (294f662; fixed upstream 1fae2d7) & $M_\pi$ / $M_\pi$ ($M_K$ with \code{-{}-compat-target 294f662})\\
\finding{U-A2} & upstream 59ccdf6 stops at the first event of a kaon run with the tree (66 values, 65 names) & 66 columns in both modes\\
\finding{D-037 F7} & the drift after a decay uses the new slopes & track ends at the drift end / Fortran\\
\finding{D-037 F4} & a particle that ranges out in the collimator continues (slow muons from kaon decays) & stopped / Fortran\\
\finding{D-061 F1} & HRS transport lengths garbled: survival weight $+\infty$ for HRS kaons & intended lengths / Fortran\\
\finding{D-061 F6} & HRS $z_\mathrm{aero}=-183$ cm does not match the HRS path end & kept in both\\
\finding{D-043 F2 data} & the $\Sigma^0$ file has stray records near threshold & kept in both\\
\finding{D-043 F7} & fixed decay model (70/30, equal-mass $p^\star$, no further decay) & kept in both\\
\finding{D-043 F8} & Saghai $\varphi$ never set (0); $\theta_{pq}$ used in the interference terms & kept in both\\
\finding{D-043 F9} & coherent production with the nucleon fit and the photon--nucleus $W$; coherent $\Lambda$ uses the $\Sigma$ branches & kept in both\\
\finding{D-043 F10} & $M_p$ in the kinematics of the factorized model for every final state & kept in both\\
\finding{D-043 F12} & $P^*_h$ message for the central event of $\Sigma$ and nuclear runs & kept in both\\
\finding{D-043 F4} & quadruple precision in the Saghai sums & reproduced\\
\finding{D-043 F14} & the relative momentum of the pair is never set for H and coherent production & not an output\\
\finding{D-043 F16} & every stop exits with status 0 & exit status 3\\
\finding{P10.1 kaon files} & tables read from the working directory & read from \file{data/kaon}, checked\\
\end{simcfindings}

\begin{simcfindingbox}
\textbf{Final-nucleus mass of coherent production.} For \code{which\_kaon} $\ge10$ the input
\code{targ\%mrec\_amu} (the hypernucleus) is first replaced by $M_A-M_p$, because it differs from that ``$A-1$'' guess
by about $M_Y$ \fsrc{dbase.f:286-292}, and then by $M_A-M_N+M_Y$ \fsrc{dbase.f:432-439}: the input value is never used,
so a hypernucleus binding energy cannot be given. The reference run \code{kaon\_he3coh\_k2c\_rad\_dec1\_A} prints both
replacement messages. The Fortran report of the kaon references describes it as the analogue of the pion coherent-mass
finding, but FORTRAN\_DIFFERENCES.md has no kaon entry. Kept in both modes.
\end{simcfindingbox}

\begin{simcrefs}
\simcref{koltenuk} D.~M.~Koltenuk, Ph.D.\ thesis (named at \fsrc{physics\_kaon.f:176-177}; the factorized model; not
  compared).
\simcref{dt} D.~Drechsel and L.~Tiator, J.\ Phys.\ G \textbf{18}, 449 (1992) (CGLN response functions; not cited by the
  code, used here only to compare the form of \eqr{eq:saghai}).
\end{simcrefs}
\end{document}
